\medskip

Soit $k$ un nombre réel.

On considère la suite $\left(u_n\right)$ définie par son premier terme $u_0$ et pour tout entier naturel $n$,

\[u_{n+1} = ku_n\left(1- u_n\right).\]





\bigskip

\textbf{Partie 1}

\medskip

Dans cette partie, $k = 1,9$ et $u_0 = 0,1$; on a donc,  pour tout entier naturel $n$,\: $u_{n+1} = 1,9u_n\left(1 - u_n\right)$.



\begin{enumerate}
\item On considère la fonction $f$ définie sur [0~;~1] par $f(x) = 1,9x(1 - x)$.
	\begin{enumerate}
		\item On étudie les variations de $f$ sur l'intervalle [0~;~1].
		
$f'(x)=1,9 \left (1-x\right ) + 1,9x\times(-1) = 1,9 - 1,9x -1,9x = 1,9\left (1-2x\right )$		
		
\begin{center}
\renewcommand{\arraystretch}{1.4}
\def\esp{\hspace*{1.5cm}}
$\begin{array}{|c | *{5}{c} |} 
\hline
x  & 0 & \esp & \frac{1}{2} & \esp  & 1 \\
\hline
f'(x) =1,9\left (1-2x\right )&  & \pmb{+} &  \vline\hspace{-2.7pt}{0} & \pmb{-} &\\
\hline
\end{array}$
\renewcommand{\arraystretch}{1}
\end{center}	

$f(0)=0$; $f\left (\frac{1}{2}\right) = 0,475$ et $f(1)=0$

On établit le tableau des variations de $f$ sur$[0\;; 1]$:

\begin{center}
{\renewcommand{\arraystretch}{1.3}
\psset{nodesep=3pt,arrowsize=2pt 3}
\def\esp{\hspace*{2.5cm}}
\def\hauteur{0pt}
$\begin{array}{|c|l*4{c}|}
\hline
x & 0  & \esp & \frac{1}{2} & \esp & 1 \\ 
\hline
f'(x) &  & \pmb{+} & \vline\hspace{-2.7pt}0 & \pmb{-} & \\ 
\hline
 & & & \Rnode{max}{0,475}  & & \\
f(x) &  & &  & & \rule{0pt}{\hauteur} \\ 
 &  \Rnode{min1}{0} &   &  &  &   \Rnode{min2}{0} \rule{0pt}{\hauteur}
 \ncline{->}{min1}{max} \ncline{->}{max}{min2} 
 \\
\hline
\end{array} $
}
\end{center}		
		
		\item On peut en déduire que si $x \in [0~;~1]$ alors $f(x) \in [0~;~0,475]$ donc $f(x) \in [0~;~1]$.
	\end{enumerate}	

\item Ci-dessous sont représentés les premiers termes de la suite $\left(u_n\right)$ construits à partir de la courbe $\mathcal{C}_f$ de la fonction $f$ et de la droite $D$ d'équation $y = x$.



\medskip


\psset{unit=10cm,arrowsize=2pt 3,comma=true}
\def\xmin{-0.3} \def\xmax{1.05}
\def\ymin{-0.2} \def\ymax{0.65} 
\begin{pspicture*}(\xmin,\ymin)(\xmax,\ymax)


\psgrid[unit=1cm,gridlabels=0,subgriddiv=1,gridcolor=gray](0,0)(1,1)

\psaxes[linewidth=1.5pt,Dx=0.1,Dy=0.1,labelFontSize=\scriptstyle]{->}(0,0)(0,0)(\xmax,\ymax)

\uput[dl](0,0){O}

\def\f{1.9 x mul  1 x sub mul}

\psplot[plotpoints=2000,linewidth=1.25pt,linecolor=blue]{0}{1}{\f}

\uput[10](0.8,0.3){\blue $\mathscr{C}_f$}

\psplot[plotpoints=1000,linecolor=red!80]{0}{1}{x}
\uput[d]{45}(0.65,0.65){\red $D : y=x$}

\psFixpoint[linecolor=red,ArrowInside=-<]{0.1}{\f}{3}
\psline[linestyle=dashed](0.1,0.171)(0,0.171)\uput[l](0,0.171){\small $u_1 = f\left(u_0\right)$}
\uput[d](0.1,-0.075){\small $u_0$}
\psline[linestyle=dashed](0.171,0.171)(0.171,0)\uput[d](0.171,-0.075){\small $u_1$}
\psline[linestyle=dashed](0.269,0.269)(0.269,0)\uput[d](0.269,-0.075){\small $u_2$}
\psline[linestyle=dashed](0.374,0.374)(0.374,0)\uput[d](0.374,-0.075){\small $u_3$}
\psline[linecolor=red](0.374,0.374)(0.374,0.445)
\end{pspicture*}


La suite $(u_n)$ semble croissante et semble converger vers l'abscisse du point d'intersection de $\mathcal{C}_f$ et $D$.

\item
	\begin{enumerate}
		\item
		
Soit $\mathcal{P}_n$ la propriété: $0 \leqslant u_n \leqslant u_{n+1} \leqslant \frac12.$

\begin{list}{\textbullet}{}
\item \textbf{Initialisation}

Pour $n=0$, $u_n=u_0=0,1$ et $u_{n+1}=1,9 u_0 \left (1-u_0\right ) = 1,9\times 0,1 \times 0,9 = 0,171$; donc $0 \leqslant u_0 \leqslant u_1 \leqslant \frac{1}{2}$.

La propriété est vraie au rang 0.

\item \textbf{Hérédité}

On suppose la propriété vraie au rang $n$, c'est-à-dire $0 \leqslant u_n \leqslant u_{n+1} \leqslant \frac12$. C'est l'hypothèse de récurrence.

La fonction $f$ est croissante sur $\left [0\;; \frac{1}{2}\right ]$ donc 
$f(0) \leqslant f(u_n) \leqslant f(u_{n+1}) \leqslant f\left (\frac12\right )$.

$f(0)=0$, $f(u_n)=u_{n+1}$, $f(u_{n+1}) = u_{n+2}$ et $f\left (\frac{1}{2}\right )=0,475 \leqslant \frac{1}{2}$

On en déduit que: $0 \leqslant u_{n+1} \leqslant u_{n+2} \leqslant \frac12$.

Donc la propriété est vraie au rang $n+1$.


\item \textbf{Conclusion}

La propriété est vraie au rang 0, et elle est héréditaire; donc, d'après le principe de récurrence, elle est vraie pour tout entier naturel $n$.
\end{list}

On a donc démontré que, pour tout entier naturel $n$, on a: $0 \leqslant u_n \leqslant u_{n+1} \leqslant \frac12$.

		\item 
\begin{list}{\textbullet}{De la propriété précédente, on tire:}
\item pour tout $n$, $u_{n}\leqslant u_{n+1}$ donc la suite $(u_n)$ est croissante;
\item pour tout $n$, $u_n\leqslant \frac{1}{2}$ donc la suite $(u_n)$ est majorée par $\frac{1}{2}$.
\end{list}

La suite $\left(u_n\right)$ est croissante et majorée donc, d'après le théorème de la convergence monotone, elle est convergente.

		\item 
On appelle $\ell$ la limite de la suite $(u_n)$. 

Comme pour tout$n$, $f\left(u_{n}\right) = u_{n+1}$ et que la fonction $f$ est continue, la limite $\ell$ vérifie $f(\ell) = \ell$.

On résout l'équation $f(x)=x$:

$\begin{gathered}f(x)=x 
\iff 1,9x(1-x)= x
\iff 1,9x(1-x)-x =0\\
\phantom{f(x)= x}
\iff x\left (1,9 - 1,9x -1\right )=0
\iff x\left (0,9 - 1,9x\right )=0\\
\phantom{f(x)=x}
\iff x=0 \text{ ou } 0,9-1,9x=0
\iff x=0 \text{ ou } x=\frac{0,9}{1,9}\end{gathered}$

La suite $(u_n)$ est croissante donc $\ell \geqslant u_0$ soit $\ell \geqslant 0,1$.

La seule solution possible est donc $\ell = \frac{0,9}{1,9}=\frac{9}{19}\approx 0,474$.
	\end{enumerate}
\end{enumerate}

\bigskip

\textbf{Partie 2}

\medskip

 Dans cette partie, $k= \frac12$  et $u_0 = \frac14$, donc, pour tout $n$,\: $u_{n+1} = \frac12 u_n \left(1 - u_n\right)$.

On admet que pour tout entier naturel $n :\: 0 \leqslant u_n \leqslant \left(\frac12\right)^n$.



\begin{enumerate}
\item 
$0<\dfrac{1}{2}<1$ donc la suite géométrique $(v_n)$ définie pour tout $n$ par $v_n=\left (\dfrac{1}{2} \right )^n$ est convergente vers 0.

Pour tout $n$, $ 0 \leqslant u_n \leqslant \left(\dfrac12\right)^n$ et $\ds\lim_{n\to +\infty} 0 = \ds\lim_{n\to +\infty} \left (\frac{1}{2}\right )^n=0$ donc, d'après le théorème des gendarmes, la suite $(u_n)$ est convergente et a pour limite 0.

\item On considère la fonction Python \texttt{algo(p)} où \texttt{p} désigne un entier naturel non nul:

\begin{center}
\begin{tabular}{|l|}\hline
\texttt{def algo(p) :}\\
\quad  \texttt{u =1/4}\\
\quad \texttt{n = 0}\\
\quad \texttt{while u > 10**(-p):}\\
\quad \qquad \texttt{u = 1/2*u*(1 - u)}\\
\quad \qquad \texttt{n = n+1}\\
\quad \texttt{return(n)}\\ \hline
\end{tabular}
\end{center}

La boucle s'arrête quand \texttt{u} est inférieur ou égal à \texttt{10**(-p)}, c'est-à-dire pour la première valeur de $n$ vérifiant $u_n \leqslant 10^{-p}$.
Comme $\ds\lim_{n\to +\infty} u_n=0$, il y a une première valeur $n_0$ à partir de laquelle $u_n \leqslant 10^{-p}$ pour tout $n\geqslant n_0$.
La boucle \texttt{while} ne tourne donc pas indéfiniment, ce qui permet à la commande \texttt{algo(p)} de renvoyer une valeur.
\end{enumerate}

\bigskip
